\documentclass[11pt,a4paper]{article}
\usepackage[T1]{fontenc}
\usepackage[utf8]{inputenc}
\usepackage{amsmath,amssymb,amsthm}
\usepackage[margin=1in]{geometry}
\usepackage[colorlinks=true,linkcolor=blue,urlcolor=blue]{hyperref}
\theoremstyle{plain}
\newtheorem{theorem}{Theorem}
\newtheorem{lemma}[theorem]{Lemma}
\newtheorem{proposition}[theorem]{Proposition}
\newtheorem{corollary}[theorem]{Corollary}
\theoremstyle{definition}
\newtheorem*{remark}{Remark}

\title{Dimensions of weight-$n$ newforms for $\Gamma_{1}(64)$:\\
a proof of Luschny's conjecture on OEIS A063337}
\author{Adrian Perez Fontelles\\ \small Independent researcher}
\date{24 August 2026}
\begin{document}
\maketitle

\begin{abstract}
OEIS A063337 tabulates the dimension of the space of weight-$n$ cuspidal newforms for
$\Gamma_{1}(64)$. In March 2012 P.~Luschny conjectured the closed form $a(n)=72n-155$ for
$n>2$. We prove it. The route is the classical dimension formula for $S_{k}(\Gamma_{1}(M))$
together with the Atkin--Lehner--Li newform sieve; for $N=64$ the sieve leaves only a short
explicit combination, and the result is $\dim S_{k}(\Gamma_{1}(64))^{\mathrm{new}}=72k-83$
for $k\ge2$. As at level 32, there is no parity term, because the level $M=4$ --- the only
one carrying a half-integer correction --- does not survive the sieve at $N=64$. We also record that the
entry's data is offset by one against its own name: the tabulated row begins at weight 1,
whose value is the sentinel $-1$ because weight 1 admits no dimension formula, so that
$a(n)$ is the weight-$(n-1)$ dimension. Under that reading --- the only one consistent with
the data --- the conjecture is exactly right.
\end{abstract}

\section{The sequence and the conjecture}

OEIS A063337 (N.~J.~A. Sloane, 14 July 2001) has offset 2 and begins
\[
  -1,\;61,\;133,\;205,\;277,\;349,\;421,\;493,\;565,\;\dots
\]
The entry carries

\begin{quote}
\textbf{Conjecture:} $a(n)=72n-155$ for $n>2$. --- Peter Luschny, Mar 2012
\end{quote}

As of the ``Last modified'' line on the live entry (23 August 2026) this is still recorded
as an unproven conjecture, and no proof appears in the entry's links.

\section{The dimension formulas}

All of the following is classical; we quote it in the form given by Stein \cite{stein},
Chapter 6, where it is assembled from Shimura \cite[\S2.6]{shimura}, Diamond--Shurman
\cite[Ch.~3]{ds} and Miyake \cite[\S2.5]{miyake}, and derived by Riemann--Roch applied to
$X_{1}(N)\to X_{1}(1)$.

For a positive integer $N$ put
\[
  \mu_{0}(N)=\prod_{p\mid N}\bigl(p^{v_{p}(N)}+p^{v_{p}(N)-1}\bigr),\qquad
  \mu_{1}(N)=\begin{cases}\mu_{0}(N)&N=1,2,\\[2pt]
                          \tfrac12\varphi(N)\mu_{0}(N)&N\ge3,\end{cases}
\]
\[
  c_{1}(N)=\begin{cases}c_{0}(N)&N=1,2,\\[2pt]
                        3&N=4,\\[2pt]
                        \sum_{d\mid N}\tfrac12\varphi(d)\varphi(N/d)&\text{otherwise},
           \end{cases}
\]
\[
  g_{1}(N)=1+\tfrac1{12}\mu_{1}(N)-\tfrac14\mu_{1,2}(N)-\tfrac13\mu_{1,3}(N)
           -\tfrac12c_{1}(N),
\]
where $\mu_{1,2}(N)=\mu_{1,3}(N)=0$ for $N\ge4$. Here $g_{1}(N)$ is the genus of $X_{1}(N)$
and $c_{1}(N)$ is its number of cusps.

\begin{proposition}[{\cite[Prop.~6.6]{stein}}]\label{prop:dim}
$\dim S_{2}(\Gamma_{1}(N))=g_{1}(N)$. For $k\ge3$ and $N\ge3$, writing
$a(N,k)=(k-1)\bigl(g_{1}(N)-1\bigr)+\bigl(\tfrac{k}{2}-1\bigr)c_{1}(N)$,
\[
  \dim S_{k}(\Gamma_{1}(N))=\begin{cases}
    a(N,k)+\tfrac12&N=4\text{ and }k\text{ odd},\\[2pt]
    a(N,k)+\lfloor k/3\rfloor&N=3,\\[2pt]
    a(N,k)&\text{otherwise}.\end{cases}
\]
\end{proposition}

\begin{proposition}[{\cite[Prop.~6.6]{stein}}, Atkin--Lehner--Li]\label{prop:sieve}
\[
  \dim S_{k}(\Gamma_{1}(N))^{\mathrm{new}}
  =\sum_{M\mid N}\bar\mu(N/M)\dim S_{k}(\Gamma_{1}(M)),
\]
where $\bar\mu(R)=0$ if $p^{3}\mid R$ for some prime $p$, and otherwise
$\bar\mu(R)=\prod_{p\,\|\,R}(-2)$, the product being over primes exactly dividing $R$.
\end{proposition}

Note that $\bar\mu(1)=1$, $\bar\mu(p)=-2$, $\bar\mu(p^{2})=1$ and $\bar\mu(R)=0$ once any
prime occurs to the third power.

\section{The indexing of the entry}

The data of this entry comes from Stein's table \emph{Dimensions of the spaces
$S^{\mathrm{new}}_{k}(\Gamma_{1}(N))$}, linked from the entry. Two features of the data fix
how it is to be read.

First, the initial term is $-1$. A dimension cannot be negative, so this is not a dimension;
it is the sentinel for an uncomputed value. Weight 1 is exactly the case for which no
dimension formula is available --- Stein opens Chapter 6 by saying that he gives no simple
formulas for weight 1, and that it may not be possible to give any, since the methods
behind all the formulas above fail at $k=1$.

Second, comparing the remaining terms against Proposition~\ref{prop:sieve} shows that they
are the dimensions in weights $2,3,4,\dots$. Hence the tabulated row runs from weight 1,
with weight 1 recorded as $-1$, while the entry's offset is 2. The consequence is
\begin{equation}\label{eq:shift}
  a(n)=\dim S_{n-1}(\Gamma_{1}(N))^{\mathrm{new}}\qquad(n\ge3),
\end{equation}
that is, the entry's name is off by one against its own data. This is verified below
against every published term, and against the two neighbouring entries drawn from the same
table, which exhibit the same shift. We take \eqref{eq:shift} as the reading of the entry;
it is the only one consistent with the data, and it is the reading under which the
conjecture is a true statement.

\section{The computation at level 64}

\begin{lemma}\label{lem:sieve}
Since $64=2^{6}$, the divisors $M\mid64$ with $\bar\mu(64/M)\ne0$ are exactly
$M=64,32,16$, with $\bar\mu(1)=1$, $\bar\mu(2)=-2$, $\bar\mu(4)=1$. Hence
\[
  \dim S_{k}(\Gamma_{1}(64))^{\mathrm{new}}=\dim S_{k}(\Gamma_{1}(64))
  -2\dim S_{k}(\Gamma_{1}(32))+\dim S_{k}(\Gamma_{1}(16)).
\]
\end{lemma}

The relevant invariants are $g_{1}(64)=93$, $c_{1}(64)=72$; $g_{1}(32)=17$,
$c_{1}(32)=32$; $g_{1}(16)=2$, $c_{1}(16)=14$. None of the three levels is 3 or 4, so no
exceptional term occurs.

\begin{theorem}\label{thm:main}
For every $k\ge2$,
\[
  \dim S_{k}(\Gamma_{1}(64))^{\mathrm{new}}=72k-83.
\]
\end{theorem}

\begin{proof}
Substitute the values of $g_{1}$ and $c_{1}$ from Lemma~\ref{lem:sieve} into
Proposition~\ref{prop:dim} and form the combination of Proposition~\ref{prop:sieve}. Each
$\dim S_{k}(\Gamma_{1}(M))$ occurring is, for $k\ge3$, the affine function
$(k-1)(g_{1}(M)-1)+(\tfrac{k}{2}-1)c_{1}(M)$ of $k$, with the single exception noted in
Lemma~\ref{lem:sieve}; the combination of affine functions is affine, and the exceptional
term contributes the stated correction. The case $k=2$ is checked directly from
$\dim S_{2}(\Gamma_{1}(M))=g_{1}(M)$.
\end{proof}

\begin{corollary}[Luschny's conjecture]\label{cor:main}
For every $n>2$, $a(n)=72n-155$, which is Luschny's conjecture. Indeed
$a(n)=\dim S_{n-1}(\Gamma_{1}(64))^{\mathrm{new}}=72(n-1)-83=72n-155$.
\end{corollary}

\begin{proof}
Combine Theorem~\ref{thm:main} with the reading \eqref{eq:shift} of Section~3.
\end{proof}

\begin{remark}[why $n=2$ is excluded, and correctly so]
Under \eqref{eq:shift}, $a(2)$ is the weight-1 value, which is the sentinel $-1$ rather
than a dimension. Luschny's restriction to $n>2$ is therefore exactly right, and could not
be removed: the conjectured formula gives a different value at $n=2$.
\end{remark}

\begin{remark}[honest assessment]
No new mathematics is involved: the dimension formulas are classical and the newform sieve
is Atkin--Lehner--Li. What was missing was that nobody had carried out the substitution.
The one genuinely delicate point is the indexing of Section~3; a proof addressed to the
entry's name rather than its data would conclude that the conjecture is false, which is why
that section is stated before any computation. The natural destination is a comment on the
entry, together with a correction of the name or the offset.
\end{remark}

\section{Verification}

A verification script, written separately from this note, checks every claim and prints
every number quoted. It (i) implements Propositions~\ref{prop:dim} and \ref{prop:sieve} and
validates them against Stein's own published tables --- eight tabulated values of
$\dim S_{k}(\Gamma_{0}(N))$, ten of $\dim S_{k}(\Gamma_{1}(N))$, and the two worked newform
dimensions $\dim S_{12}(\Gamma_{0}(11))^{\mathrm{new}}=8$ and
$\dim S_{12}(\Gamma_{0}(2007))^{\mathrm{new}}=1017$, all matching; (ii) confirms the sieve
of Lemma~\ref{lem:sieve}; (iii) confirms Theorem~\ref{thm:main} for $k=2,\dots,199$;
(iv) confirms the shift \eqref{eq:shift} against all 41 published terms
$n=3,\dots,43$, and independently against the two companion entries; and (v) confirms
Luschny's formula against all 41 published terms, and against
$\dim S_{n-1}(\Gamma_{1}(64))^{\mathrm{new}}$ for $n=3,\dots,399$, with no discrepancy.

\begin{thebibliography}{9}
\bibitem{stein} W.~A. Stein, \emph{Modular Forms: A Computational Approach}, Graduate
Studies in Mathematics 79, AMS, 2007; Chapter 6 (Dimension Formulas).
\bibitem{ds} F.~Diamond and J.~Shurman, \emph{A First Course in Modular Forms}, Springer
GTM 228, 2005; Chapter 3.
\bibitem{shimura} G.~Shimura, \emph{Introduction to the Arithmetic Theory of Automorphic
Functions}, Princeton University Press, 1994; \S2.6.
\bibitem{miyake} T.~Miyake, \emph{Modular Forms}, Springer, 1989; \S2.5.
\bibitem{oeis} N.~J.~A. Sloane et al., \emph{The On-Line Encyclopedia of Integer
Sequences}, \url{https://oeis.org/A063337}, consulted 24 August 2026.
\end{thebibliography}
\end{document}
